\documentclass[11pt]{article}
\usepackage{graphicx}
\usepackage{amsmath}
\usepackage{amssymb}
\usepackage{booktabs}
\usepackage{hyperref}

\title{High-Fidelity Simulation and Detection of Hallucinations Across GPT Model Generations}
\author{Anonymous Author}
\date{}

\begin{document}
\maketitle

\begin{abstract}
Large language models (LLMs) can produce fluent but incorrect outputs, a phenomenon often called \emph{hallucination}. This paper presents a reproducible framework for (i) generating challenging evaluation prompts, (ii) labeling model outputs with an interpretable hallucination detector, and (iii) analyzing failure patterns across model families. Our evaluation mixes two fixed, human-curated prompt sets (50 factual questions and 50 riddles) with a third set of 50 \emph{auto-generated} questions that is freshly sampled each run to probe mild distribution shift. We compare GPT-3.5 Turbo, GPT-4, GPT-4o, and GPT-5 using the same harness and report both hallucination rates and abstention (no-output) behavior. Averaged over two runs, hallucination rates decrease substantially with newer models: GPT-3.5 Turbo averages 23.7\%, GPT-4 15.3\%, GPT-4o 8.0\%, and GPT-5 4.0\%. Notably, GPT-5 shows near-zero hallucinations on riddles, but frequently abstains (empty output; 20\% overall), indicating a coverage--risk trade-off. A human audit of 449 labeled examples shows the detector achieves high precision (0.92) and recall (0.97) for hallucination identification, supporting its use for comparative analysis and debugging.
\end{abstract}

\section{Introduction}
Hallucination in generative AI refers to outputs that contain unsupported or false claims presented as if they were true. This behavior poses reliability risks in applications that require grounded information, and it remains one of the most visible failure modes of modern LLMs. Benchmarks such as TruthfulQA were designed to elicit plausible but incorrect answers, demonstrating that models can mimic human misconceptions and that greater raw capability does not automatically yield truthfulness \cite{truthfulqa}. Alignment and instruction-following post-training, including reinforcement learning from human feedback (RLHF), have been shown to improve reliability and reduce misleading behavior \cite{instructgpt,gpt4report}. 

Alongside mitigation, \emph{detection} of hallucinations remains challenging. Hallucinations vary in form (incorrect entities, wrong numbers, subtle contradictions, or confidently stated false explanations). Surveys emphasize that robust evaluation often benefits from combining multiple signals and maintaining interpretability so failures can be debugged \cite{ji_survey}. This paper contributes a practical, reproducible harness and an interpretable detector for reference-based settings where an expected answer is available.

\paragraph{Contributions.} Our contributions are:
\begin{enumerate}
    \item A reproducible evaluation harness that combines fixed curated prompts with re-sampled auto-generated prompts for stress-testing under mild distribution shift.
    \item An interpretable hallucination detector that aggregates lexical overlap, numeric consistency checks, and contradiction/unavailability cues, with optional extensions for embedding similarity and bidirectional NLI contradiction signals.
    \item A cross-model study of hallucination and abstention behavior across GPT-3.5 Turbo, GPT-4, GPT-4o, and GPT-5, using a shared evaluation protocol.
    \item Pattern analysis highlighting topic- and keyword-level triggers for hallucinations in GPT-3.5 and showing that simple length-based surface features have weak correlation with hallucination.
\end{enumerate}

\section{Related Work}
TruthfulQA measures model truthfulness under adversarial questioning and highlights how models may reproduce common false beliefs rather than grounded facts \cite{truthfulqa}. Post-training alignment methods, including instruction tuning and RLHF, can reduce hallucination and improve reliability in practice \cite{instructgpt}. The GPT-4 technical report discusses extensive alignment and safety work and reports improvements on measures related to factuality and reliability \cite{gpt4report}.

For hallucination detection, prior work in summarization and dialogue includes approaches that measure factual consistency between generated text and a reference source. FactCC, for example, targets factual consistency in abstractive summarization \cite{factcc}. More broadly, surveys catalog a spectrum of strategies: lexical matching, semantic similarity, entailment/contradiction modeling, and retrieval-based verification \cite{ji_survey}. Our work focuses on \emph{reference-based question answering} (QA): each prompt has an expected answer (gold or silver-standard). This setting allows interpretable detection rules and controlled comparisons across model families.

\section{Methodology}
\subsection{Problem Setup}
For each prompt $x$, we assume a reference answer $y^\star$ (gold or silver-standard). Given a model response $\hat{y}$, the detector outputs a binary label $\texttt{hallucinated} \in \{0,1\}$. We treat hallucination as the presence of at least one \emph{contradictory} or \emph{inconsistent} factual claim relative to $y^\star$. We explicitly distinguish hallucination from \emph{abstention}: responses that provide no answer (or provide no content) contain no false factual claim but also do not satisfy the task. In the released harness, if post-processing yields an empty response, the run records an abstention and sets \texttt{hallucinated=False} for that item. In analysis, we report both hallucination rate and no-output (abstention) rate.

\subsection{Detector Signals}
The detector is implemented as a set of interpretable signals designed to catch common hallucination modes while remaining debuggable.

\paragraph{Lexical overlap and containment.}
We compute lexical similarity using the Jaccard similarity between content-token sets extracted from $\hat{y}$ and $y^\star$:
\begin{equation}
J(A,B)=\frac{|A\cap B|}{|A\cup B|}.
\end{equation}
We also perform normalized substring containment and token subsequence containment checks for short factual answers (e.g., named entities or terms). Lexical overlap is inexpensive and informative, but can miss paraphrases and can be fooled by responses that share keywords while changing a crucial detail.

\paragraph{Numeric consistency and conflicts.}
We extract numeric strings (including years and percentages) and apply:
\begin{itemize}
    \item Exact-match constraints for year/identifier-like numbers.
    \item Tolerance-based matching for general numeric quantities:
    \begin{equation}
    |\hat{r}-r^\star|\le \delta_{\text{abs}}+\delta_{\text{rel}}|r^\star|,
    \end{equation}
    where $\delta_{\text{rel}}=5\times 10^{-3}$ and $\delta_{\text{abs}}=5\times 10^{-2}$.
    \item A disjunction heuristic that flags answers offering mutually exclusive numeric alternatives (e.g., ``1491 or 1492'') as likely incorrect.
\end{itemize}

\paragraph{Contradiction and unavailability cues.}
We flag explicit contradiction patterns such as negating an entity that exists in the reference (e.g., ``no such X exists'') and unavailability cues (e.g., ``there is no information about Y'') when the reference provides a concrete answer.

\paragraph{Optional semantic similarity and NLI.}
The codebase supports optional extensions: sentence-embedding cosine similarity via Sentence-BERT \cite{sentencebert} and bidirectional natural language inference (NLI) contradiction signals trained on MNLI \cite{mnli} (with RoBERTa as a standard backbone \cite{roberta}). These extensions can be enabled when computational resources allow; the core detector remains usable without them.

\subsection{Decision Rule}
We use a \emph{high-precision} decision rule: a response is labeled hallucinated when there is strong evidence of contradiction or numeric inconsistency (or an atomic entity mismatch under strict conditions). Otherwise, it is labeled non-hallucinated and the detector records diagnostic metadata (e.g., overlap scores and which checks fired). This design prioritizes avoiding false positives; a limitation is that off-topic answers that do not directly contradict the reference may be labeled non-hallucinated. We address this in analysis by separately reporting abstentions and by validating the detector against human labels (Section~\ref{sec:audit}).

\section{Experimental Setup}
\subsection{Prompt Sets}
Each evaluation run contains 150 prompts, split evenly:
\begin{enumerate}
    \item Factual (50): Curated short-answer questions across diverse topics (e.g., astronomy, biology, law, finance). References emphasize the minimum correct fact.
    \item Riddle (50): Classic riddles and logic puzzles with canonical solutions.
    \item Auto-generated (50): A newly sampled set of niche questions produced by an LLM at run time. We generate 10 topical seed domains with 5 questions each (50 total) and then generate a silver-standard reference answer for each question via a separate answer-generation step (implemented in the released harness). Because this set is re-sampled each run, it probes mild distribution shift and reduces overfitting to a fixed list.
\end{enumerate}
Across runs, the factual and riddle sets are fixed; the auto-generated set is independently re-sampled.

\subsection{Models and Generation Parameters}
We evaluate four GPT model endpoints via the OpenAI API: GPT-3.5 Turbo, GPT-4, GPT-4o, and GPT-5. All models receive the same high-level instruction wrapper requesting concise answers. The harness uses default sampling parameters in code: GPT-3.5 Turbo uses temperature $1.0$ by default; GPT-4 and GPT-4o use temperature $0.8$ by default; GPT-5 is queried via the Responses API and may omit sampling parameters depending on model capabilities (see the released code for exact request fields). We run two evaluation rounds. For factual and riddle prompts, the two rounds measure run-to-run variability on a fixed set. For auto-generated prompts, the second round is an independent re-sample, capturing both model stochasticity and prompt-set difficulty variation.

\subsection{Human Audit of Detector Labels}
\label{sec:audit}
To validate the detector, we construct a stratified audit set by sampling labeled examples across categories and models, then manually annotate whether hallucination is present. Table~\ref{tab:audit} shows results on 449 annotated examples: precision 0.92 and recall 0.97. This supports using the detector as a reliable high-precision labeler for comparative analysis and debugging.

\begin{table}[t]
\centering
\begin{tabular}{lrrrr}
\toprule
 & TP & FP & FN & TN \\
\midrule
Detector vs.\ Human (449) & 60 & 5 & 2 & 382 \\
\bottomrule
\end{tabular}
\caption{Human audit of hallucination labels. Precision $=60/(60+5)=0.92$, recall $=60/(60+2)=0.97$, accuracy $=(60+382)/449=0.98$.}
\label{tab:audit}
\end{table}

\section{Results}
\subsection{Hallucination and Abstention Rates}
Table~\ref{tab:hallu_rates} reports hallucination rates by category, averaged over two runs (with min--max ranges across the runs in parentheses). Because the auto-generated set is re-sampled each run, its range reflects both model stochasticity and dataset difficulty variation.

\begin{table}[t]
\centering
\begin{tabular}{lcccc}
\toprule
Model & Factual (50) & Riddle (50) & Auto-gen (50) & Overall (150) \\
\midrule
GPT-3.5 Turbo & 19.0 (18--20) & 29.0 (28--30) & 23.0 (22--24) & 23.7 (23.3--24.0) \\
GPT-4 & 7.0 (6--8) & 25.0 (24--26) & 14.0 (8--20) & 15.3 (12.7--18.0) \\
GPT-4o & 7.0 (6--8) & 16.0 (16--16) & 1.0 (0--2) & 8.0 (8.0--8.0) \\
GPT-5 & 8.0 (6--10) & 0.0 (0--0) & 4.0 (0--8) & 4.0 (2.0--6.0) \\
\bottomrule
\end{tabular}
\caption{Hallucination rates (\%) by model and prompt category. Values are the mean of two runs; ranges show min--max across runs.}
\label{tab:hallu_rates}
\end{table}

We also quantify no-output (abstention) rates---cases where the model response is empty/blank---in Table~\ref{tab:abstain}. GPT-5 shows substantial abstention (20\% overall) and especially high abstention on riddles (mean 32\%), while other models rarely abstain in this harness.

\begin{table}[t]
\centering
\begin{tabular}{lcccc}
\toprule
Model & Factual & Riddle & Auto-gen & Overall \\
\midrule
GPT-3.5 Turbo & 0.0 & 0.0 & 0.0 & 0.0 \\
GPT-4 & 0.0 & 0.0 & 0.0 & 0.0 \\
GPT-4o & 0.0 & 0.0 & 0.0 & 0.0 \\
GPT-5 & 7.0 (6--8) & 32.0 (30--34) & 21.0 (18--24) & 20.0 (20--20) \\
\bottomrule
\end{tabular}
\caption{No-output (abstention) rates (\%) where the model response is empty/blank. Values are means with min--max ranges over two runs.}
\label{tab:abstain}
\end{table}

\subsection{Prompt-Level Patterns}
For GPT-3.5 Turbo, hallucinations concentrate in specific topical templates and prompt cues. Figure~\ref{fig:template} shows pronounced template sensitivity in one run: some templates exhibit elevated rates, but template counts are uneven; extreme values (e.g., 100\%) can occur for templates with very small sample sizes in the 150-prompt set. Overall, math/science-leaning templates show higher error rates than many others in this run. Figure~\ref{fig:keywords} conditions hallucination rates on the presence of TF--IDF keywords selected from hallucinated prompts, highlighting prompt cues that co-occur with errors (e.g., riddle-related tokens such as ``legs''). Finally, Figure~\ref{fig:heatmap} shows that prompt length and response length have only weak correlation with hallucination ($|r| \le 0.16$), indicating that content and domain characteristics dominate over simple surface features.

\begin{figure}[t]
\centering
\includegraphics[width=0.95\linewidth]{figures/hallucinations_by_template.png}
\caption{GPT-3.5 Turbo hallucination rate by topical template (run A; all 150 prompts).}
\label{fig:template}
\end{figure}

\begin{figure}[t]
\centering
\includegraphics[width=0.95\linewidth]{figures/hallucinations_by_keywords.png}
\caption{GPT-3.5 Turbo hallucination rate conditioned on TF--IDF keywords (selected from hallucinated prompts; run A).}
\label{fig:keywords}
\end{figure}

\begin{figure}[t]
\centering
\includegraphics[width=0.95\linewidth]{figures/feature_correlation_heatmap.png}
\caption{Feature correlation heatmap (GPT-3.5 Turbo, run A). Correlations between hallucination and prompt/response length are small in magnitude.}
\label{fig:heatmap}
\end{figure}

\subsection{Run-to-Run Variability on Fixed Prompt Sets}
Because factual and riddle prompts are identical across the two runs, we measure how often the hallucination label flips on these 100 fixed prompts. GPT-3.5 Turbo flips on 10\% of prompts, GPT-4 flips on 3\%, GPT-4o flips on 5\%, and GPT-5 flips on 4\%. Variability is highest for riddles in GPT-3.5 and GPT-4, reinforcing that riddle-style prompts amplify sensitivity to stochastic generation even under consistent instruction formatting.

\section{Discussion}
The results show consistent improvements in hallucination rates across newer GPT model families. However, aggregate hallucination rate alone does not capture the full reliability story: GPT-5 exhibits a substantial abstention (no-output) mode, especially on riddles. This suggests a coverage--risk trade-off: abstaining reduces the chance of stating a falsehood but also reduces answer completeness. In settings where incorrect answers are costly, abstentions can be preferable; in other settings, they may be frustrating and call for fallback mechanisms (retrieval, re-asking, or escalation to a human).

The prompt-level analyses highlight actionable debugging insights: GPT-3.5 shows strong topic sensitivity (e.g., finance and certain science templates), and specific lexical cues correlate with hallucination in riddles. Meanwhile, the weak correlation between hallucination and simple length features suggests that detectors should focus on semantic content and consistency rather than shallow heuristics. Finally, the human audit indicates that interpretable, rule-based signals can yield high-precision hallucination labeling suitable for tracking model trends and targeting mitigation.

\paragraph{Limitations.} The auto-generated references are silver-standard and may contain occasional errors; reference quality bounds evaluation quality. The high-precision detector can miss off-topic answers that do not directly contradict the reference, so hallucination rates should be interpreted as a lower bound on overall incorrectness. Additionally, abstention is operationalized here as empty output; future work should distinguish explicit refusals from API failures and quantify abstention in richer categories (e.g., hedged answers, partial answers, and safe refusals).

\section{Conclusion}
We presented a reproducible evaluation and analysis framework for hallucinations across GPT model families, including a modular, interpretable detector and a hybrid prompt suite combining fixed and re-sampled generated questions. A human audit supports the detector's reliability, and the empirical results show a substantial reduction in hallucination rates across model generations. We also observe an important behavioral shift in GPT-5: near-zero hallucinations on riddles accompanied by substantial abstention via no-output responses. Future work can incorporate retrieval-based grounding for the abstained cases, broaden the audited label set, and test the same protocol on additional model families and open-source baselines.

\begin{thebibliography}{99}

\bibitem{truthfulqa}
S.~Lin, J.~Hilton, and O.~Evans,
``TruthfulQA: Measuring How Models Mimic Human Falsehoods,''
in \emph{Proc. ACL}, 2022.

\bibitem{ji_survey}
Z.~Ji, N.~Lee, R.~Frieske, F.~Yu, D.~Su, Y.~Xu, E.~I. Fung, and A.~Kumar,
``Survey of Hallucination in Natural Language Generation,''
\emph{ACM Computing Surveys}, 55(12), 2023.

\bibitem{instructgpt}
L.~Ouyang, J.~Wu, X.~Jiang, D.~Almeida, C.~Wainwright, P.~Mishkin, et~al.,
``Training Language Models to Follow Instructions with Human Feedback,''
in \emph{Advances in Neural Information Processing Systems (NeurIPS)}, 2022.

\bibitem{gpt4report}
OpenAI,
``GPT-4 Technical Report,''
\emph{arXiv:2303.08774}, 2023.

\bibitem{factcc}
W.~Kry{\'s}ci{\'n}ski, B.~McCann, C.~Xiong, and R.~Socher,
``Evaluating the Factual Consistency of Abstractive Text Summarization,''
in \emph{Proc. EMNLP}, 2019.

\bibitem{sentencebert}
N.~Reimers and I.~Gurevych,
``Sentence-BERT: Sentence Embeddings using Siamese BERT-Networks,''
in \emph{Proc. EMNLP-IJCNLP}, 2019.

\bibitem{mnli}
A.~Williams, N.~Nangia, and S.~Bowman,
``A Broad-Coverage Challenge Corpus for Sentence Understanding through Inference,''
in \emph{Proc. NAACL}, 2018.

\bibitem{roberta}
Y.~Liu, M.~Ott, N.~Goyal, J.~Du, M.~Joshi, D.~Chen, et~al.,
``RoBERTa: A Robustly Optimized BERT Pretraining Approach,''
\emph{arXiv:1907.11692}, 2019.

\end{thebibliography}

\end{document}